\documentclass[11pt]{article}

\usepackage{amssymb,amsmath,amsthm}
\usepackage{authblk}
\usepackage{multirow}
\usepackage{blindtext}
\usepackage{amsfonts}
%\usepackage{amsmath,amssymb,xcolor}
\usepackage{amsmath,amssymb,pxfonts,xcolor}
\usepackage{slashed}
\usepackage{graphicx}
%\usepackage{natbib}
\usepackage{comment}
\usepackage{ulem}
\usepackage{color}

\oddsidemargin 0.0in
\textwidth 6.3in
\setcounter{topnumber}{5}
\setcounter{bottomnumber}{5}
\setcounter{totalnumber}{10}
\newtheorem{theorem}{Theorem}
\newtheorem{acknowledgement}[theorem]{Acknowledgement}
\newtheorem{algorithm}[theorem]{Algorithm}
\newtheorem{axiom}[theorem]{Axiom}
\newtheorem{case}[theorem]{Case}
\newtheorem{claim}[theorem]{Claim}
\newtheorem{conclusion}[theorem]{Conclusion}
\newtheorem{condition}[theorem]{Condition}
\newtheorem{conjecture}[theorem]{Conjecture}
\newtheorem{corollary}[theorem]{Corollary}
\newtheorem{criterion}[theorem]{Criterion}
\newtheorem{definition}[theorem]{Definition}
%\newtheorem{example}[theorem]{Example}
\newtheorem{exercise}[theorem]{Exercise}
\newtheorem{lemma}[theorem]{Lemma}
\newtheorem{notation}[theorem]{Notation}
\newtheorem{problem}[theorem]{Problem}
\newtheorem{proposition}[theorem]{Proposition}
\newtheorem{remark}[theorem]{Remark}
\newtheorem{solution}[theorem]{Solution}
\newtheorem{summary}[theorem]{Summary}


%\numberwithin{equation}{section}

\title{A lower bound for the volatility swap in the lognormal SABR model}
\author{Elisa Al\`os \quad Frido Rolloos \quad Kenichiro Shiraya\thanks{Kenichiro Shiraya is supported by CARF.}}


\begin{document}
\maketitle

\abstract{In the short time to maturity limit it is proved that for the conditionally lognormal SABR model the zero vanna implied volatility is a lower bound for the volatility swap strike. The result is valid for all values of the correlation parameter and is a sharper lower bound than the at-the-money implied volatility for correlation less than or equal to zero.}


\section{Introduction}
Liquid instruments to manage volatility risk include variance and volatility swaps. The former is substantially easier to price than the latter because volatility swaps are more exposed to model risk. From a risk management and pricing standpoint it is therefore desirable to have robust apprroximations and bounds for volatility swap prices.

\smallskip
One robust approximation is given in the important paper of Carr and Lee \cite{CL1}. By making use of an immunization technique, and assuming that the skew is generated by a stochastic volatility (SV) model, Carr and Lee were able to give an approximate price for the volatility swap with an error of $\mathrm O(\rho^2)$. Even though their result is approximate when correlation is nonzero, it remains an important result as the approximation is independent of the specific form of the SV model.

\smallskip
A more recent approximation for the volatility swap is given by Rolloos and Arslan \cite{RA}. In the aforementioned paper it is shown that the zero vanna implied volatility (ZVIV) approximates the volatility swap price. The error between the exact volatility swap price and the ZVIV was made precise by Al\`os et al. \cite{ARS} by making use of techniques from Malliavin calculus. As the volatility swap price is model-dependent, the error between the ZVIV and the exact price will depend on the particular SV model.

\smallskip
In the aforementioned paper by Al\`os et al. numerical simulations indicate that in addition to approximating volatility swaps the ZVIV is also a lower bound for the volatility swap. In this paper it is proved that for the lognormal SABR model (henceforth SABR model) this is indeed the case for all values of correlation in the short time to maturity limit. To the best of our knowledge this has not been proved yet in the literature on derivatives pricing. A proof for general SV models is a subject of future research.

\smallskip
The structure of the note is as follows. In Section 2 the working assumptions are stated and relevant notation introduced. Section 3 treats the zero correlation case. It is proved that both the at-the-money implied volatility (ATMI) and ZVIV are lower bounds for the volatility swap in the short maturity limit in the SABR model. By making use of the result that the ZVIV is a better approximation than the ATMI it follows that the ZVIV is a sharper lower bound. The nonzero correlation case is discussed in Section 4, wherein it is proved that in the SABR model the ZVIV is a lower bound for the volatility swap price for any value of correlation. Furthermore, it is shown that for non-positive correlation it is a sharper lower bound than the ATMI. Section 5 concludes.



\section{Assumptions and notation}

Consider a stochastic volatility model for the log-price of a stock
under a risk-neutral probability measure $P$:
\begin{align}
\label{themodel}
X_{t} &=X_0-\frac{1}{2}\int_{0}^{t}{\sigma _{s}^{2}}ds+\int_{0}^{t}\sigma
_{s}\left( \rho dW_{s}+\sqrt{1-\rho ^{2}}dB_{s}\right) \\
\label{SABR}
\sigma_t &= \sigma_0 + \int_0^t \alpha \sigma_s \, dW_s
\end{align}
Here, $X_0$ is the current log-price, $%
W$ and $B$ are independent standard Brownian motions defined on a complete probability
space  $(\Omega ,\mathcal{G},P)$, and $\sigma $ is a square-integrable and
right-continuous stochastic process adapted to the filtration generated by $%
W $.  The volatility of volatility $\alpha$ is a positive constant. We denote by $\mathcal{F}^{W}$ and $\mathcal{F}^{B}$
the filtrations generated by $W$ and $B$ and $\mathcal{F}:=%
\mathcal{F}^{W}\vee \mathcal{F}^{B}.$ We assume the interest rate $r$ to be zero for the sake of simplicity. The same arguments in this paper hold for $r\neq 0$.


Under the above model, the price of a European call with strike price $K$ is given by the equality
\[
V_{t}=E_{t}[(e^{X_{T}}-K)_{+}], 
\]%
where $E_{t}$ is the $\mathcal{F}_{t}-$conditional expectation with respect
to $P$ (i.e., $E_{t}[Z]=E[Z|\mathcal{F}_{t}]$). In the sequel, we make use
of the following notation:

\begin{itemize}
\item $v_t=\sqrt{\frac{Y_t}{T-t}}$, where $Y_t=\int_{t}^{T}\sigma _{u}^{2}du$.

That is, $v$
represents the future average volatility, and it is not an adapted process. 
Notice that $E_{t}\left[
v_{t}\right] $ is the fair strike of a volatility swap with maturity time $T.$

\item $BS(t,T,x,k,\sigma )$ is the price of a  European call option
under the classical Black-Scholes model with constant volatility $\sigma $,
stock price $e^x$, time to maturity $T-t,$ and strike $K=\exp
(k) $. Remember that (if $r=0$)

\[
BS(t,T,x,k,\sigma )=e^{x}N(d_1(k,\sigma ))-e^{k}N(d_2(k,\sigma )), 
\]%
where $N$ denotes the cumulative probability function of the standard normal
law and

\[
d_1\left( k,\sigma \right) :=\frac{x-k}{\sigma \sqrt{T-t}}+
\frac{\sigma }{2}\sqrt{T-t},  \hspace{0.4cm} d_2\left( k,\sigma \right) :=\frac{x-k}{\sigma \sqrt{T-t}}-
\frac{\sigma }{2}\sqrt{T-t}.
\]%
For the sake of simplicity we make use of the notation $BS(k,\sigma):=BS(t,T,X_t,k,\sigma )$.
\item  The inverse function $BS^{-1}(t,T,x,k, \cdot)$ of the Black-Scholes formula with respect to the volatility parameter is defined as
\[
BS(t,T,x,k, BS^{-1}( t,T,x,k,\lambda) )=\lambda,
\]
for all $\lambda>0$.
For the sake of simplicity, we denote $BS^{-1}(k,\lambda)\ :=BS^{-1}(t,T,X_{t},k,\lambda)$.

\item For any fixed $t,T,X_{t},k,$ we define the implied volatility $I(t,T,X_{t},k) $ as the quantity such that 
\[
BS( t,T,X_{t},k,I( t,T,X_{t},k) ) =V_{t}.
\]
Notice that $I(t,T,X_t,k)=BS^{-1}( t,T,X_t,k,V_t)$.

\item $k_t^*$ denotes the ATM strike and $I(t,T,X_t,k_t^*)$ the corresponding ATMI

\item $\hat{k_t}$ is the {\it zero vanna strik}e at time $t$. That is, the strike such that 
$$d_2(\hat{k}_t,I(t,T,X_t,\hat{k}_t))=0.$$ Moreover, we will refer to $I(t,T,X_t,\hat{k}_t)$ as the {\it zero vanna implied volatility}. We recall that (Black-Scholes) vanna is the partial derivative of the Black-Scholes delta with respect to implied volatility, which is directly proportional to $d_2$. The zero vanna strike is therefore the strike where the Black-Scholes vanna of a vanilla option is zero.

\item $
\Lambda _{r}:=E_{r}\left[ BS\left( t,T,X_t,\hat{k}_t,v_{t}\right) \right]$.

\item $
\Theta_{r}(k):=BS^{-1}(k,\Lambda_r)$. Notice that $\Theta_{t}(\hat{k}_t)=I(t,X_t,\hat{k}_t,V_t)$ and $\Theta_{T}(\hat{k}_t)=v_t$.

\item $G(t,T,x,k,\sigma ):=( \frac{\partial ^{2}}{\partial x^{2}}-\frac{%
\partial}{\partial x}) BS(t,T,x,k,\sigma )$.
\item $H(t,T,x,k,\sigma ):=( \frac{\partial ^{3}}{\partial x^{3}}-\frac{%
\partial ^{2}}{\partial x^{2}}) BS(t,T,x,k,\sigma )$.
%\item $K(t,T,x,k,\sigma ):=( \frac{\partial ^{3}}{\partial x^{3}}-\frac{\partial ^{2}}{\partial x^{2}}) ( \frac{\partial ^{2}}{\partial x^{2}}-\frac{\partial }{\partial x})BS(t,T,x,k,\sigma )$

\end{itemize}


In the remaining of this paper $\mathbb{D}_{W}^{1,2}$ denotes the domain of the
Malliavin derivative operator $D^{W}$ with respect to the Brownian motion $%
W. $ We also consider the
iterated derivatives $D^{n,W}$ , for $n>1,$ whose domains will be denoted by 
$\mathbb{D}_{W}^{n,2}$. We will use the notation $\mathbb{L}_{W}^{n,2}=$\ $%
L^{2}(\left[ 0,T\right] ;\mathbb{D}_{W}^{n,2})$. Notice that the SABR variance $\sigma^2$ given by (\ref{SABR}) is a process in $\mathbb{L}_{W}^{2,2}$ and, moreover
\begin{eqnarray}
D_r\sigma_u^2 &=& 2\alpha \sigma_u^2 1_{[0,u)}(r)\\
D_s D_r \sigma_u^2 &=& 2 D_s\sigma_u D_r\sigma_u + 2\sigma_u D_s D_r \sigma_u\nonumber\\
&=& 2 \alpha^2 \sigma_u^2 1_{[0,u)}(r \vee s) + 2 \alpha^2 \sigma_u^2 1_{[0,u)}(r \vee s)\nonumber\\
&=& 4 \alpha^2 \sigma_u^2 1_{[0,u)}(r \vee s).
\end{eqnarray}

\section{The uncorrelated case}

\subsection{ATM implied volatility}
\begin{theorem} Consider the SABR model defined as in (\ref{themodel}) and (\ref{SABR}). Then, for zero correlation, the ATMI is a lower bound for the volatility swap in the short time to maturity limit.
\begin{equation}
\lim_{T \to t} \, I( t,T,X_t,k^*_t) \leq \lim_{T \to t} \,  E_{t}[v_{t}]
\end{equation}
\begin{proof}
This follows immediately from Theorem 3.2 of Al\`os and Shiraya \cite{AS},
\begin{eqnarray}
\lim_{T\rightarrow t}\frac{I\left( t,T,X_t,k^*_t\right) -E_{t}\left[v_{t}\right] }{(T-t)^{2}}
= -\frac{1}{32\sigma _{t}}\lim_{T\rightarrow t}
\frac{1}{(T-t)^{3}}E_{t}\left[\int_{t}^{T}\left(E_{r}\left[\int_{r}^{T}D_{r}^{W}\sigma_{s}^{2}ds\right]\right) ^{2}dr\right] \le 0.
\end{eqnarray}
\end{proof}
\end{theorem}


\subsection{Zero-vanna implied volatility}

\begin{theorem}Consider the SABR model defined as in (\ref{themodel}) and (\ref{SABR}). Then, for zero correlation, the ZVIV is a lower bound for the volatility swap in the short time to maturity limit.
\begin{equation}
\lim_{T \to t} \, I( t,T,X_t,\hat{k}_t) \leq \lim_{T \to t} \,  E_{t}[v_{t}]
\end{equation}
\begin{proof}
In the uncorrelated case, since (B.3) of Al\`os, Rolloos and Shiraya,
\begin{eqnarray}
\label{primerpas}
I\left( t,T,X_t,\hat{k}_t\right) &=& E_{t}\left[ v_{t}\right] -  E_t\Bigg[\int_{t}^{T}\left( BS^{-1}\left( \hat{k}_t, \Lambda_r\right)\right) ^{\prime \prime\prime} (D^-A)_r U_r dr\Bigg]\nonumber\\
&&- \frac12 E_t\Bigg[
\int_{t}^{T}\left( BS^{-1}\left( \hat{k}_t, \Lambda_r\right)\right) ^{(iv)} A_r U_r^2 dr \Bigg].
\label{I-E[v]}\nonumber\\
&=& 
E_{t}\left[ v_{t}\right] -\left( BS^{-1}\left( \hat{k}_t, \Lambda_t\right)\right) ^{\prime \prime\prime} E_t\Bigg[\int_{t}^{T}(D^-A)_r U_r dr\Bigg]\nonumber\\
&&+ \frac12 \left( BS^{-1}\left( \hat{k}_t, \Lambda_t\right)\right) ^{(iv)} E_t\Bigg[
\int_{t}^{T} A_r U_r^2 dr \Bigg]\nonumber\\
&&+o((T-t)^{4H+1})\nonumber\\
&=& E_{t}\left[ v_{t}\right] + T_1 + T_2 +o((T-t)^{4H+1})
\end{eqnarray}
where 
\begin{eqnarray}
A_r&:=&\frac12 \int_{r}^{T} U_{s}^{2}ds,\\
(D^-A)_r&:=&\frac{1}{2}\int^T_r D_r^W U^2_s ds,\\
U_{r}&:=&E_{r}\left[D_r^W\left(BS(t,T,X_t,\hat{k}_t,v_t)\right)\right]\nonumber\\
&=&E_r\left[\frac{\partial BS}{\partial \sigma}(t,T,X_t,\hat{k}_t,v_t)\frac{1}{2v_t(T-t)}\int_r^T D_r^W\sigma_s^2 ds\right],
\end{eqnarray}
and $\Theta_{r}(k):=BS^{-1}(k,\Lambda_r)$.

Here,
\begin{eqnarray}
\label{terceraderivada}
\left( BS^{-1}\left( \hat{k}_t, \Lambda_t\right)\right) ^{\prime \prime\prime}
&=&(2\pi)^{\frac32}\exp\left(-3X_t+\frac{3}{2}(\Theta_t(\hat{k}_t))^2(T-t)\right)(T-t)^{-\frac12}+o\left((T-t)^{-\frac{1}{2}}\right)\\
\label{derivada4}
\left( BS^{-1}\left( \hat{k}_t, \Lambda_t\right)\right) ^{(iv)}&=&-\frac{3(2\pi)^{2}}{\Theta_t(\hat{k}_t)}\exp\left(-4X_t+2(\Theta_t(\hat{k}_t))^2(T-t)\right)(T-t)^{-1} + o\left((T-t)^{-1}\right).
\end{eqnarray}
Thus, $(BS^{-1})'''>0$ and $(BS^{-1})^{(iv)}<0$ in the short time limit.

On the other hand, 
\begin{eqnarray}
\lefteqn{E_t\Bigg[\int_{t}^{T} A_r U_r^2 dr \Bigg]}\nonumber\\
&=&\frac{1}{2}E_t\Bigg[\int_{t}^{T} \left(\int_r^TU_s^2ds\right) U_r^2 dr \Bigg]\nonumber\\
&=&\frac{1}{4}E_t\Bigg[ \left(\int_t^TU_r^2dr\right)^2  \Bigg]\nonumber\\
&=&\frac{1}{4}E_t\Bigg[ \left(\int_t^T\left(    E_r\left[\frac{\partial BS}{\partial \sigma}(t,T,X_t,\hat{k}_t,v_t)\frac{1}{2v_t(T-t)}\int_r^TD_r^W\sigma_s^2 dr\right]
    \right)^2ds\right)^2  \Bigg]\nonumber\\
&>& 0
\end{eqnarray}
Thus, $T_2\le 0$.

Next,
\begin{eqnarray}
\label{U}
U_s&=&E_s\left[\frac{\partial BS}{\partial \sigma}(t,T,X_t,\hat{k}_t,v_t)\frac{1}{2v_t(T-t)}\int_s^T D_s^W \sigma_u^2 du\right]  \nonumber\\
&=&\frac12    E_s\left[G(t,T,X_t,\hat{k}_t,v_t)\int_s^T D_s^W \sigma_u^2 du\right],
\end{eqnarray}
and
\begin{eqnarray}
\label{DU}
D_r^WU_s&=& E_s\Bigg[ \frac12 
G(t,T,X_t,\hat{k}_t,v_t)\left(\frac{d_1(\hat{k}_t,v_t)d_2(\hat{k}_t,v_t)}{2v_t^2(T-t)} - \frac{1}{2v_t^2(T-t)}\right)
\left(\int_s^T D_s^W\sigma_u^2 du\right)\left(\int_r^TD_r^W\sigma_u^2 du\right)
\nonumber\\
&&+\frac12G(t,T,X_t,\hat{k}_t,v_t)\left(\int_s^TD_r^WD_s^W\sigma_u^2 du\right)\Bigg].
\end{eqnarray}
\begin{eqnarray}
\label{DA2}
(D^-A)_r
&=&
\frac{1}{4}
\int_r^T E_s\left[\frac{e^{\hat{k}_t}N'(d_2(\hat{k}_t,v_t))}{v_t\sqrt{T-t}}\int_s^TD_s^W\sigma_u^2 du    \right] \nonumber\\
&&\times E_s \Bigg[
\frac{e^{\hat{k}_t}N'(d_2(\hat{k}_t,v_t))}{v_t\sqrt{T-t}}
\left(
\frac{-1}{2 v_t^2(T-t)}
\left(\int_s^TD_s^W\sigma_u^2 du\right)\left(\int_r^TD_r^W\sigma_u^2 du\right)+\left(\int_{r\vee s}^TD_r^WD_s^W\sigma_u^2 du\right)\right)
\Bigg] ds\nonumber\\
&&
+ O(\nu^3 (T-t)^{3H+\frac{1}{2}}),
\end{eqnarray}
Thus, we need to check
\begin{equation}
\frac{-1}{2 v_t^2(T-t)}\int_s^TD_s^W\sigma_u^2 du\int_r^TD_r^W\sigma_u^2 du+\int_s^TD_r^WD_s^W\sigma_u^2 du \ge 0,
\label{eq9}
\end{equation}
Since $\int_t^T \sigma_u^2 du \ge \int_r^T \sigma_u^2 du$, \eqref{eq9} is
\begin{eqnarray}
\lefteqn{\frac{-1}{2 \int_t^T \sigma_u^2 du}\int_s^T 2\alpha \sigma_u^2 du\int_r^T 2\alpha \sigma_u^2 du+\int_{r\vee s}^T 4\alpha^2 \sigma_u^2 du} \nonumber\\
&\ge& - \int_{r\vee s}^T 2\alpha^2 \sigma_u^2 du+\int_{r\vee s}^T 4\alpha^2 \sigma_u^2 du \nonumber\\
&\ge& 0
\end{eqnarray}
Thus, $T_1 \le 0$ and the zero-vanna iv is a lower bound for the volatility swap.
\end{proof}
\end{theorem}


\section{The correlated case}

\subsection{ATM implied volatility}

\begin{theorem}\label{correlatedATMI}
Consider the SABR model defined as in (\ref{themodel}) and (\ref{SABR}). Then, for $\rho < 0$, the ATMI is a lower bound for the volatility swap strike in the short time to maturity limit.
\begin{equation}
\lim_{T \to t} \, I(t,T,X_t,k^*_t) \leq \lim_{T \to t} \,  E_{t}[v_{t}]
\end{equation}
\begin{proof}
From Theorem 4.2 in \cite{AS},
\begin{eqnarray}
\label{lim1}
\lefteqn{\lim_{T\rightarrow t}\frac{I( t,T,X_t,{k}_t^*) -E_{t}[v_{t}]} {(T-t)}}\nonumber \\
&=& \lim_{T\rightarrow t}\frac{3\rho^{2}}{8\sigma _{t}^{3}(T-t)^{4}}
E_{t}\left[\left(\int_{t}^{T}\int_{s}^{T}D_{s}^{W}\sigma _{r}^{2}dr ds\right)^{2} \right]\nonumber\\
&&-\lim_{T\rightarrow t}\frac{\rho ^{2}}{{2\sigma _{t}^2(T-t)^{3}}}
E_{t}\left[ \int_{t}^{T} \int_{s}^{T}D_{s}^{W}\sigma_{r} \int_{r}^{T}D_{r}^{W}\sigma _{u}^{2}du dr ds\right]\nonumber\\
&&-\lim_{T\rightarrow t}\frac{\rho ^{2}}{2\sigma _{t}(T-t)^{3}}
E_{t}\left[\int_{t}^{T} \int_{s}^{T}\int_{r}^{T}D_{s}^{W}D_{r}^{W}\sigma _{u}^{2}du dr ds\right]\nonumber\\
&&+\lim_{T\rightarrow t}\frac{\rho}{4(T-t)^{2}}
E_{t}\left[\int_{t}^{T}\int_{s}^{T}D_{s}^{W}\sigma _{r}^{2}dr ds \right]\nonumber\\
&=& \lim_{T\rightarrow t}\frac{3\rho^{2}}{8\sigma _{t}^{3}(T-t)^{4}}T_1 - \lim_{T\rightarrow t}\frac{\rho ^{2}}{{2\sigma _{t}^2(T-t)^{3}}} T_2 - \lim_{T\rightarrow t}\frac{\rho ^{2}}{2\sigma _{t}(T-t)^{3}} T_3 +\lim_{T\rightarrow t}\frac{\rho}{4(T-t)^{2}} T_4.
\end{eqnarray}

Here, in the log normal stochastic volatility model, since Schwartz's inequality,
\begin{eqnarray}
T_1 
%&=& 2 \int_{t}^{T} \int_0^s E_{t}\left[ \int_{s}^{T}D_{s}^{W}\sigma _{r}^{2}dr \int_{u}^{T}D_{u}^{W}\sigma _{v}^{2}dv \right] du ds\nonumber\\
&\le& 
(T-t)\int_{t}^{T}(T-s)\int_{s}^{T}E_{t}\left[\left(D_{s}^{W}\sigma _{r}^{2}\right)^{2}\right] dr ds \nonumber\\
&=&
(T-t)\int_{t}^{T}(T-s)\int_{s}^{T}4\alpha^2 E_{t}\left[ \sigma_0^4 e^{4\alpha W_r - 2\alpha^2 r}\right] dr ds \nonumber\\
&=&
(T-t)\int_{t}^{T}(T-s)\int_{s}^{T}4\alpha^2 \sigma_0^4 e^{6 \alpha^2 r} dr ds \nonumber\\
&=&
4\alpha^2 \sigma_0^4 (T-t)\int_{t}^{T}(T-s)\frac{e^{6 \alpha^2 T} - e^{6 \alpha^2 s}}{6\alpha^2} ds \nonumber\\
&=&
4\alpha^2 \sigma_0^4  \left(\frac{e^{6 \alpha^2 T}(T-t)^3}{12\alpha^2} + \frac{e^{6 \alpha^2 t}}{36\alpha^4}(T-t)^2 - (T-t)\int_t^T \frac{e^{6 \alpha^2 s}}{36\alpha^4}ds \right)\nonumber\\
&=&
4\alpha^2 \sigma_0^4  \left(\frac{e^{6 \alpha^2 T}(T-t)^3}{12\alpha^2} + \frac{e^{6 \alpha^2 t}}{36\alpha^4}(T-t)^2 - \frac{e^{6 \alpha^2 T} - e^{6 \alpha^2 t}}{216\alpha^6}(T-t) \right)\label{t1}
\end{eqnarray}
Here, we assume $t=0$ and apply the Taylor expansion,
\begin{eqnarray}
\eqref{t1}
&=& 
4\alpha^2 \sigma_0^4  \left(\frac{1+6\alpha^2T }{12\alpha^2}T^3 + \frac{T^2}{36\alpha^4} - \frac{6\alpha^2T + \frac{1}{2}36\alpha^4 T^2 + \frac{1}{6}216\alpha^6T^3}{216\alpha^6}T\right) + O(T^5)\nonumber\\
&=& 
4\alpha^2 \sigma_0^4  \left(\frac{1}{2}T^4 - \frac{1}{6}T^4\right) + O(T^5)\nonumber\\
&=& 
\frac{4 \alpha^2 \sigma_0^4}{3}T^4+O(T^5)
\end{eqnarray}

\begin{eqnarray}
T_2 
&=& \int_{t}^{T} \int_{s}^{T} \int_{r}^{T} E_{t}\left[ D_{s}^{W}\sigma_{r}  D_{r}^{W}\sigma _{u}^{2} \right]du dr ds\nonumber\\
&=& \int_{t}^{T} \int_{s}^{T} \int_{r}^{T} 2\alpha^2 \sigma_0^3  E_{t}\left[ e^{2\alpha(W_u - W_r) - 2\alpha^2(u-r)}\right]E_{t}\left[ e^{3\alpha W_r - \frac12 9\alpha^2r}\right]e^{\alpha^2u + 2\alpha^2r}du dr ds\nonumber\\
&=& 2\alpha^2 \sigma_0^3 \int_{t}^{T} \int_{s}^{T} \int_{r}^{T} e^{\alpha^2u + 2\alpha^2r}du dr ds\nonumber\\
&=& 2\alpha^2 \sigma_0^3 \int_{t}^{T} \int_{s}^{T} \frac{e^{\alpha^2T + 2\alpha^2r} - e^{3\alpha^2r}}{\alpha^2} dr ds\nonumber\\
&=& 2\alpha^2 \sigma_0^3 \int_{t}^{T} \frac{e^{3\alpha^2T} - e^{\alpha^2T + 2\alpha^2s}}{2\alpha^4} - \frac{e^{3\alpha^2T} - e^{3\alpha^2s}}{3\alpha^4} dr ds\nonumber\\
&=& 2\alpha^2 \sigma_0^3 \left( \frac{e^{3\alpha^2T}(T-t)}{2\alpha^4} - \frac{e^{3\alpha^2T} - e^{\alpha^2T + 2\alpha^2t}}{4\alpha^6} - \frac{e^{3\alpha^2T}(T-t)}{3\alpha^4} + \frac{e^{3\alpha^2T} - e^{3\alpha^2t}}{9\alpha^6} \right)\label{t2}
\end{eqnarray}
Here, we assume $t=0$ and apply the Taylor expansion,
\begin{eqnarray}
\eqref{t2}
%&=& 
%\lim_{T\to 0} 2\alpha^2 \sigma_0^3  \left(\frac{e^{3\alpha^2T}T}{6\alpha^4} - \frac{2\alpha^2 e^{\alpha^2T}T}{4\alpha^6} + \frac{3\alpha^2 T}{9\alpha^6} +O(T^2)\right)\nonumber\\
&=& 
2\alpha^2 \sigma_0^3  \Bigg(\frac{(1+3\alpha^2T+\frac{1}{2}9\alpha^4T^2)T}{6\alpha^4} - \frac{(2\alpha^2T + \frac{1}{2}4\alpha^4T^2 + \frac{1}{6}8\alpha^6T^3)(1 + \alpha^2T + \frac{1}{2}\alpha^4T^2)}{4\alpha^6}\nonumber\\
&& + \frac{3\alpha^2 T+ \frac{1}{2}9\alpha^4T^2+\frac{1}{6}27\alpha^6T^3}{9\alpha^6}\Bigg) +O(T^4)\nonumber\\
&=&
\frac{\alpha^2 \sigma_0^3}{3}T^3+O(T^4)
\end{eqnarray}

\begin{eqnarray}
T_3 
&=& \int_{t}^{T} \int_{s}^{T}\int_{r}^{T} E_{t}\left[ 4 \alpha^2\sigma_0 e^{2\alpha W - \alpha^2 u}1_{[0,u)}(r \vee s)\right]du dr ds\nonumber\\
&=& \int_{t}^{T} \int_{s}^{T}\int_{r}^{T} 4 \alpha^2\ \sigma_0^2 e^{\alpha^2 u}1_{[0,u)}(r \vee s)du dr ds\nonumber\\
&=& 4 \alpha^2\ \sigma_0^2 \int_{t}^{T} \int_{s}^{T} \frac{e^{\alpha^2 T} - e^{\alpha^2 r}}{\alpha^2} dr ds\nonumber\\
&=& 4 \alpha^2\ \sigma_0^2 \int_{t}^{T} \frac{e^{\alpha^2 T}(T-s)}{\alpha^2} - \frac{e^{\alpha^2 T} - e^{\alpha^2 s}}{\alpha^4} ds\nonumber\\
&=& 4 \alpha^2\ \sigma_0^2 \left( \frac{e^{\alpha^2 T}(T-t)^2}{2\alpha^2} - \frac{e^{\alpha^2 T}(T-t)}{\alpha^4} + \frac{e^{\alpha^2 T} - e^{\alpha^2 t}}{\alpha^6} \right)\label{t3}
\end{eqnarray}
Here, we assume $t=0$ and apply the Taylor expansion,
\begin{eqnarray}
\eqref{t3}
&=& 
4\alpha^2 \sigma_0^2  \left(\frac{1+\alpha^2T}{2\alpha^2}T^2 - \frac{1+\alpha^2T + \frac{1}{2}\alpha^4T^2}{\alpha^4}T + \frac{\alpha^2T + \frac{1}{2}\alpha^4T^2 + \frac{1}{6}\alpha^6T^3}{\alpha^6}\right)+O(T^4)\nonumber\\
&=& 
\frac{2 \alpha^2\sigma_0^2}{3}T^3+O(T^4)
\end{eqnarray}
\begin{eqnarray}
T_4
&=& 
\int_{t}^{T}\int_{s}^{T}E_{t}\left[D_{s}^{W}\sigma _{r}^{2}\right] dr ds \nonumber\\
&=&
\int_{t}^{T}\int_{s}^{T}2\alpha E_{t}\left[ \sigma_0^2 e^{2\alpha W_r - \alpha^2 r}\right] dr ds \nonumber\\
&=&
\int_{t}^{T}\int_{s}^{T}2\alpha \sigma_0^2 e^{\alpha^2 r} dr ds \nonumber\\
&=&
2\alpha \sigma_0^2 \int_{t}^{T} \frac{e^{\alpha^2 T} - e^{\alpha^2 s}}{\alpha^2} ds \nonumber\\
&=&
2\alpha \sigma_0^2  \left(\frac{e^{\alpha^2 T}(T-t)}{\alpha^2} - \frac{e^{\alpha^2 T} - e^{\alpha^2 t}}{\alpha^4} \right)\label{t4}
\end{eqnarray}
Here, we assume $t=0$ and apply the Taylor expansion,
\begin{eqnarray}
\eqref{t4}
&=& 
2\alpha \sigma_0^2  \left(\frac{1+\alpha^2 T}{\alpha^2}T - \frac{\alpha^2T + \frac{1}{2}\alpha^4 T^2}{\alpha^4}\right) + O(T^3)\nonumber\\
&=& 
\alpha \sigma_0^2  \left(T^2 - \frac{1}{2}T^2\right) + O(T^3)\nonumber\\
&=& 
\frac{\alpha \sigma_0^2}{2}T^2+O(T^3)
\end{eqnarray}

Thus,
\begin{eqnarray}
\lefteqn{\lim_{T\rightarrow 0}\frac{I( 0,T,X_0,k_0^*) -E_{0}[v_{0}]} {T}}\nonumber \\
&\le& 
\lim_{T\rightarrow 0}\frac{3\rho^{2}}{8\sigma _{0}^{3}T^{4}} \frac{4 \alpha^2 \sigma_0^4}{3}T^4
- \lim_{T\rightarrow 0}\frac{\rho ^{2}}{{2\sigma _{0}^2T^{3}}}\frac{\alpha^2 \sigma_0^3}{3}T^3
- \lim_{T\rightarrow 0}\frac{\rho ^{2}}{2\sigma _{0}T^{3}} \frac{2 \alpha^2\sigma_0^2}{3}T^3
+ \lim_{T\rightarrow 0}\frac{\rho}{4T^{2}} \frac{\alpha \sigma_0^2}{2}T^2 + \lim_{T\to 0}o(T)\nonumber\\
&=& \frac{\rho \alpha \sigma_0^2}{8}
\end{eqnarray}

%%%%%%%%%%%%%%%%%%%%%%%%%%%%%%%%%%%%%%%%%%%%%%%%%%%%%%%%%%%%%%%%%%%%%%%%%%%%%%%%%%%%%%%%%%%%%%%%%%%%%%%%%%%
\begin{comment}
By Proposition 4.1 in \cite{AS},
\begin{eqnarray}
I\left( t,T,X_t,k_t\right) &=&I^{0}\left( t,T,X_t,k_t\right)
\nonumber \\
&&+\frac{\rho }{2}E_t \left[\int_{t}^{T}\left( BS^{-1}\right) ^{\prime }\left( k_t,\Lambda _{s}\right) H(s,T,X_{s},k_t,v_{s})\Phi _{s}ds\right],
\label{impliedrel}
\end{eqnarray}%
where $I^{0}\left( t,T,X_t,k_t\right) $ denotes the implied
volatility in the uncorrelated case $\rho =0$,
\begin{equation}
\Lambda _{s}:=E_{t}\left[ BS(t,T,X_t,k_t,v_{t})\right] +\frac{\rho }{2}%
E_{t}\left[\int_{t}^{s} H(r,T,X_{r},k_t,v_{r})\Phi _{r}dr\right], 
\end{equation}
\begin{equation}
\Phi _{t}:=\sigma _{t}\int_{t}^{T}D_{t}^{W}\sigma _{r}^{2}dr
\end{equation}
\begin{equation}
(BS^{-1})'(k,\Lambda_r) = \frac{1}{e^{X_t}N'(D_+(k_t,\Gamma_r))\sqrt{T-t}},
\end{equation}
\begin{eqnarray}
H(r,T,X_{r},k_t,v_{r})=\frac{N'({\color{red}{d_1}})}{v_r \sqrt{T-r}}e^{X_r}(1-\frac{{\color{red}{d_1}}}{v_r\sqrt{T-r}}).
\end{eqnarray}
Here, in the ATMI case, since $x-{\color{red}k_t^*}=0$, 
\begin{eqnarray}
1-\frac{{\color{red}{d_1}}}{v_r\sqrt{T-r}}
&=& 1 - \frac12\nonumber\\
&>& 0
\end{eqnarray}
\end{comment}
%%%%%%%%%%%%%%%%%%%%%%%%%%%%%%%%%%%%%%%%%%%%%%%%%%%%%%%%%%%%%%%%%%%%%%%%%%%%%%%%%%%%%%%%%%%%%%%%%%%%%%%%%%%
Thus, ATMI is a lower bound of the volatility swap if $\rho<0$.
\end{proof}
\end{theorem}

\begin{comment}
On the other hand, in the zero-vanna iv case, since $x-k = \frac{I(k_2)^2(T-t)}{2}$,
\begin{eqnarray}
1-\frac{\color{red}{d_1}}{v_r\sqrt{T-r}}
&=& 1 - \left(\frac{I(k_2)^2(T-t)}{2 v_t^2(T-r)} + \frac{1}{2}\right)
\end{eqnarray}
Thus, it is not clear whether zero-vanna iv is a lower bound or not...

\end{comment}
\subsection{Zero-vanna implied volatility}
%%%%%%%%%%%%%%%%%%%%%%%%%%%%%%%%%%%%%%%%%%%%%%%%%%%%%%%%%%%%%%%%%%%%%%%%%%%%%%%%%%%%%%%%%%%%%%%%%%%%%%%%%%%
\begin{theorem} Consider the SABR model defined as in (\ref{themodel}) and (\ref{SABR}). Then, for all values of correlation $\rho$, the ZVIV is a lower bound for the volatility swap strike in the short time to maturity limit.
\begin{equation}
\lim_{T \to t} \, I(t,T,X_t,\hat{k}_t) \leq \lim_{T \to t} \,  E_{t}[v_{t}]
\end{equation}
\begin{proof}
From Theorem 4 in \cite{ARS},
\begin{eqnarray}
\label{lim1}
\lefteqn{\lim_{T\rightarrow t}\frac{I( t,T,X_t,\hat{k}_t) -E_{t}[v_{t}]} {(T-t)}}\nonumber \\
&=& \lim_{T\rightarrow t}\frac{3\rho^{2}}{8\sigma _{t}^{3}(T-t)^{4}}
E_{t}\left[\left(\int_{t}^{T}\int_{s}^{T}D_{s}^{W}\sigma _{r}^{2}dr ds\right)^{2} \right]\nonumber\\
&&-\lim_{T\rightarrow t}\frac{\rho ^{2}}{{2\sigma _{t}^2(T-t)^{3}}}
E_{t}\left[ \int_{t}^{T} \int_{s}^{T}D_{s}^{W}\sigma_{r} \int_{r}^{T}D_{r}^{W}\sigma _{u}^{2}du dr ds\right]\nonumber\\
&&-\lim_{T\rightarrow t}\frac{\rho ^{2}}{2\sigma _{t}(T-t)^{3}}
E_{t}\left[\int_{t}^{T} \int_{s}^{T}\int_{r}^{T}D_{s}^{W}D_{r}^{W}\sigma _{u}^{2}du dr ds\right]\nonumber\\
&=& \lim_{T\rightarrow t}\frac{3\rho^{2}}{8\sigma _{t}^{3}(T-t)^{4}}T_1 - \lim_{T\rightarrow t}\frac{\rho ^{2}}{{2\sigma _{t}^2(T-t)^{3}}} T_2 - \lim_{T\rightarrow t}\frac{\rho ^{2}}{2\sigma _{t}(T-t)^{3}} T_3.
\end{eqnarray}
%%%%%%%%%%%%%%%%%%%%%%%%%%%%%%%%%%%%%%%%%%%%%%%%%%%%%%%%%%%%%%%%%%%%%%%%%%%%%%%%%%%%%%%%%%%%%%%%%%%%%%%%%%%
\begin{comment}
Here, in the log normal stochastic volatility model, since Schwartz's inequality,
\begin{eqnarray}
T_1 
%&=& 2 \int_{t}^{T} \int_0^s E_{t}\left[ \int_{s}^{T}D_{s}^{W}\sigma _{r}^{2}dr \int_{u}^{T}D_{u}^{W}\sigma _{v}^{2}dv \right] du ds\nonumber\\
&\le& 
(T-t)\int_{t}^{T}(T-s)\int_{s}^{T}E_{t}\left[\left(D_{s}^{W}\sigma _{r}^{2}\right)^{2}\right] dr ds \nonumber\\
&=&
(T-t)\int_{t}^{T}(T-s)\int_{s}^{T}4\alpha^2 E_{t}\left[ \sigma_0^4 e^{4\alpha W_r - 2\alpha^2 r}\right] dr ds \nonumber\\
&=&
(T-t)\int_{t}^{T}(T-s)\int_{s}^{T}4\alpha^2 \sigma_0^4 e^{6 \alpha^2 r} dr ds \nonumber\\
&=&
4\alpha^2 \sigma_0^4 (T-t)\int_{t}^{T}(T-s)\frac{e^{6 \alpha^2 T} - e^{6 \alpha^2 s}}{6\alpha^2} ds \nonumber\\
&=&
4\alpha^2 \sigma_0^4  \left(\frac{e^{6 \alpha^2 T}(T-t)^3}{12\alpha^2} + \frac{e^{6 \alpha^2 t}}{36\alpha^4}(T-t)^2 - (T-t)\int_t^T \frac{e^{6 \alpha^2 s}}{36\alpha^4}ds \right)\nonumber\\
&=&
4\alpha^2 \sigma_0^4  \left(\frac{e^{6 \alpha^2 T}(T-t)^3}{12\alpha^2} + \frac{e^{6 \alpha^2 t}}{36\alpha^4}(T-t)^2 - \frac{e^{6 \alpha^2 T} - e^{6 \alpha^2 t}}{216\alpha^6}(T-t) \right)\label{t1}
\end{eqnarray}
Here, we assume $t=0$ and apply the Taylor expansion,
\begin{eqnarray}
\eqref{t1}
&=& 
4\alpha^2 \sigma_0^4  \left(\frac{1+6\alpha^2T }{12\alpha^2}T^3 + \frac{T^2}{36\alpha^4} - \frac{6\alpha^2T + \frac{1}{2}36\alpha^4 T^2 + \frac{1}{6}216\alpha^6T^3}{216\alpha^6}T\right) + O(T^5)\nonumber\\
&=& 
4\alpha^2 \sigma_0^4  \left(\frac{1}{2}T^4 - \frac{1}{6}T^4\right) + O(T^5)\nonumber\\
&=& 
\frac{4 \alpha^2 \sigma_0^4}{3}T^4+O(T^5)
\end{eqnarray}

\begin{eqnarray}
T_2 
&=& \int_{t}^{T} \int_{s}^{T} \int_{r}^{T} E_{t}\left[ D_{s}^{W}\sigma_{r}  D_{r}^{W}\sigma _{u}^{2} \right]du dr ds\nonumber\\
&=& \int_{t}^{T} \int_{s}^{T} \int_{r}^{T} 2\alpha^2 \sigma_0^3  E_{t}\left[ e^{2\alpha(W_u - W_r) - 2\alpha^2(u-r)}\right]E_{t}\left[ e^{3\alpha W_r - \frac12 9\alpha^2r}\right]e^{\alpha^2u + 2\alpha^2r}du dr ds\nonumber\\
&=& 2\alpha^2 \sigma_0^3 \int_{t}^{T} \int_{s}^{T} \int_{r}^{T} e^{\alpha^2u + 2\alpha^2r}du dr ds\nonumber\\
&=& 2\alpha^2 \sigma_0^3 \int_{t}^{T} \int_{s}^{T} \frac{e^{\alpha^2T + 2\alpha^2r} - e^{3\alpha^2r}}{\alpha^2} dr ds\nonumber\\
&=& 2\alpha^2 \sigma_0^3 \int_{t}^{T} \frac{e^{3\alpha^2T} - e^{\alpha^2T + 2\alpha^2s}}{2\alpha^4} - \frac{e^{3\alpha^2T} - e^{3\alpha^2s}}{3\alpha^4} dr ds\nonumber\\
&=& 2\alpha^2 \sigma_0^3 \left( \frac{e^{3\alpha^2T}(T-t)}{2\alpha^4} - \frac{e^{3\alpha^2T} - e^{\alpha^2T + 2\alpha^2t}}{4\alpha^6} - \frac{e^{3\alpha^2T}(T-t)}{3\alpha^4} + \frac{e^{3\alpha^2T} - e^{3\alpha^2t}}{9\alpha^6} \right)\label{t2}
\end{eqnarray}
Here, we assume $t=0$ and apply the Taylor expansion,
\begin{eqnarray}
\eqref{t2}
%&=& 
%\lim_{T\to 0} 2\alpha^2 \sigma_0^3  \left(\frac{e^{3\alpha^2T}T}{6\alpha^4} - \frac{2\alpha^2 e^{\alpha^2T}T}{4\alpha^6} + \frac{3\alpha^2 T}{9\alpha^6} +O(T^2)\right)\nonumber\\
&=& 
2\alpha^2 \sigma_0^3  \Bigg(\frac{(1+3\alpha^2T+\frac{1}{2}9\alpha^4T^2)T}{6\alpha^4} - \frac{(2\alpha^2T + \frac{1}{2}4\alpha^4T^2 + \frac{1}{6}8\alpha^6T^3)(1 + \alpha^2T + \frac{1}{2}\alpha^4T^2)}{4\alpha^6}\nonumber\\
&& + \frac{3\alpha^2 T+ \frac{1}{2}9\alpha^4T^2+\frac{1}{6}27\alpha^6T^3}{9\alpha^6}\Bigg) +O(T^4)\nonumber\\
&=&
\frac{\alpha^2 \sigma_0^3}{3}T^3+O(T^4)
\end{eqnarray}

\begin{eqnarray}
T_3 
&=& \int_{t}^{T} \int_{s}^{T}\int_{r}^{T} E_{t}\left[ 4 \alpha^2\sigma_0 e^{2\alpha W - \alpha^2 u}1_{[0,u)}(r \vee s)\right]du dr ds\nonumber\\
&=& \int_{t}^{T} \int_{s}^{T}\int_{r}^{T} 4 \alpha^2\ \sigma_0^2 e^{\alpha^2 u}1_{[0,u)}(r \vee s)du dr ds\nonumber\\
&=& 4 \alpha^2\ \sigma_0^2 \int_{t}^{T} \int_{s}^{T} \frac{e^{\alpha^2 T} - e^{\alpha^2 r}}{\alpha^2} dr ds\nonumber\\
&=& 4 \alpha^2\ \sigma_0^2 \int_{t}^{T} \frac{e^{\alpha^2 T}(T-s)}{\alpha^2} - \frac{e^{\alpha^2 T} - e^{\alpha^2 s}}{\alpha^4} ds\nonumber\\
&=& 4 \alpha^2\ \sigma_0^2 \left( \frac{e^{\alpha^2 T}(T-t)^2}{2\alpha^2} - \frac{e^{\alpha^2 T}(T-t)}{\alpha^4} + \frac{e^{\alpha^2 T} - e^{\alpha^2 t}}{\alpha^6} \right)\label{t3}
\end{eqnarray}
Here, we assume $t=0$ and apply the Taylor expansion,
\begin{eqnarray}
\eqref{t3}
&=& 
4\alpha^2 \sigma_0^2  \left(\frac{1+\alpha^2T}{2\alpha^2}T^2 - \frac{1+\alpha^2T + \frac{1}{2}\alpha^4T^2}{\alpha^4}T + \frac{\alpha^2T + \frac{1}{2}\alpha^4T^2 + \frac{1}{6}\alpha^6T^3}{\alpha^6}\right)+O(T^4)\nonumber\\
&=& 
\frac{2 \alpha^2\sigma_0^2}{3}T^3+O(T^4)
\end{eqnarray}

Thus,
\end{comment}
%%%%%%%%%%%%%%%%%%%%%%%%%%%%%%%%%%%%%%%%%%%%%%%%%%%%%%%%%%%%%%%%%%%%%%%%%%%%%%%%%%%%%%%%%%%%%%%%%%%%%%%%%%%
Since $T_1$, $T_2$, $T_3$ are the same as those in the proof of Theorem \ref{correlatedATMI}, we obtain
\begin{eqnarray}
\lefteqn{\lim_{T\rightarrow 0}\frac{I( 0,T,X_0,\hat{k}_0) -E_{0}[v_{0}]} {T}}\nonumber \\
&\le& 
\lim_{T\rightarrow 0}\frac{3\rho^{2}}{8\sigma _{0}^{3}T^{4}} \frac{4 \alpha^2 \sigma_0^4}{3}T^4
- \lim_{T\rightarrow 0}\frac{\rho ^{2}}{{2\sigma _{0}^2T^{3}}}\frac{\alpha^2 \sigma_0^3}{3}T^3
- \lim_{T\rightarrow 0}\frac{\rho ^{2}}{2\sigma _{0}T^{3}} \frac{2 \alpha^2\sigma_0^2}{3}T^3 + \lim_{T\to 0}o(T)\nonumber\\
&=&0
\end{eqnarray}
\end{proof}
\end{theorem}


\section{Conclusion}
It has been shown that in the conditionally lognormal SABR model, the ZVIV is a lower bound for the volatility swap strike for all values of correlation in the short time to maturity limit. This property, in addition to the results in \cite{ARS}, further cements the ZVIV as a more accurate model-free approximation for the volatility swap strike than the ATMI.

The results in this paper is a first step in analyzing the question for which models the ZVIV is a lower bound for the volatility swap strike. Further research into the ZVIV as a lower bound for the volatility swap strike will have to includes analyses along the dimensions of the Hurst parameter, correlation parameter and time to maturity.

\begin{thebibliography}{9}

%\bibitem{A} Al\`{o}s, E. ``A generalization of the Hull and White formula with applications to option pricing approximation,'' Finance and Stochastics 10 (3), (2006): 353-365.

%\bibitem{AL} Al\`{o}s, E. and Le\'{o}n, J. A. ``On the curvature of the smile in stochastic volatility models,'' SIAM J. Financial Math. 8, (2017): 373-399.

%\bibitem{ACTW} Alòs, E., Chatterjee, R.  Tudor, S. F. and Wang, T.H. `` Target volatility option pricing in the lognormal fractional SABR model", Quantitative Finance, 19:8, (2019): 1339-1356,

%\bibitem{ALV} Al\`{o}s, E., Le\'{o}n, J. A., and Vives, J. ``On the short-time behavior of the implied volatility for jump-diffusion models with stochastic volatility,'' Finance and Stochastics 11 (4), (2007): 571-589.

\bibitem{AS} Al\`{o}s, E., and Shiraya, K. ``Estimating the Hurst parameter from short term volatility swaps: a Malliavin calculus approach.'' Finance and Stochastics 23.2 (2019): 423-447.

\bibitem{ARS} Al\`{o}s, E., Rolloos, F., and Shiraya, K. ``On the difference between the volatility swap strike and the zero vanna implied volatility'' SIAM Journal on Financial Mathematics.

%\bibitem{BFG}Bayer, C., Friz, P. and  Gatheral, J.  ``'Pricing under rough volatility." Quantitative Finance, 16:6 (2016), pp. 887-904.

%\bibitem{bg} Bergomi, L. and Guyon, J.. The Smile in Stochastic Volatility Models (December 2, 2011). Available at SSRN: https://ssrn.com/abstract=1967470

\bibitem{CL1} Carr, P. and Lee, R. ``Robust replication of volatility derivatives,'' PRMIA award for Best Paper in Derivatives, MFA 2008 Annual Meeting, (2008).

%\bibitem{CL} Carr, P. and Lee, R. ``Volatility derivatives,'' Annu. Rev. Financ. Econ. 1 (1), (2009): 319-339.

%\bibitem{CR} Comte, F. and Renault, E.: ``Long memory in continuous-time stochastic volatility models,''  Math. Financ. 8, (1998): 291-323.

%\bibitem{DDKZ} Demeterfi, K., Derman, E., Kamal, M., and Zou, J. ``More Than You Ever Wanted To Know About Volatility Swaps,'' Goldman Sachs Quantitative Strategies Research Notes, (1999).

%\bibitem{DU2} Decreusefond, L. and Ust\"unel, A. S. . ``Stochastic analysis of the fractional Brownian motion". Potential Anal. 10, (1999), pp. 177–214.

%\bibitem{efgr} El Euch, O. Fukasawa, M., Gatheral, J. and Rosenbaum, M. ``Short-Term At-the-Money Asymptotics Under Stochastic Volatility Models,''  SIAM Journal on Financial Mathematics 10.2 (2019): 491-511.

%\bibitem{Fe} Feinstein S. ``The black-scholes formula is nearly linear in sigma for at-the-money options: Therefore implied volatilities from at-the-money options are virtually unbiased,'' Work. Pap., Fed. Reserve Bank Atlanta, (1989).

%\bibitem{F} Fukasawa, M. ``Asymptotic analysis for stochastic volatility: martingale expansion,'' Finance and Stochastics 15 (4), (2011): 635-654.

%\bibitem{F} Fukasawa, M. ``Volatility derivatives and model-free implied leverage,'' International Journal of Theoretical and Applied Finance 17 (1), (2014).

%\bibitem{FG} Friz, P., and Gatheral, J. ``Valuation of volatility derivatives as an inverse problem,'' Quantitative Finance 5 (6), (2005): 531-542.

%\bibitem{GJ} Gatheral, J. and Jaisson, T. ``Fractional volatility models. Presentation at Bloomberg seminar,''
%http://mfe.baruch.cuny.edu/wp-content/uploads/2012/09/FractionalVolatility2014BBQ.pdf, (2014).

%\bibitem{GJR} Gatheral, J., Jaisson, T. and Rosenbaum, M. ``Volatility is rough.'' Preprint, 2014. arXiv:1410.3394.

%\bibitem{N} Nualart, D. ``The Malliavin Calculus and Related Topics. Second Edition,'' Probability and its Applications. Springer-Verlag, (2006).

%\bibitem{NR} Nualart, D. and R\u{a}\c{s}canu, A. ``Differential equations driven by fractional Brownian motion,'' Collect. Math. 53 (1), (2002): 55-81.

\bibitem{RA} Rolloos, F., and Arslan, M. ``Taylor-made volatility swaps,'' Wilmott, (2017).

%\bibitem{RT} Romano, M., and Touzi, N. ``Contingent claims and market completeness in a stochastic volatility model,'' Math. Financ. 7, (1997): 399-410.

%\bibitem{W} Willard, G.A. ``Calculating prices and sensitivities for path-independent securities in multifactor models,'' J. Deriv. 5, (1997): 45-61.

\end{thebibliography}


\end{document}